\documentclass[10pt,a4paper]{amsart}
\usepackage[margin=19mm]{geometry}
\usepackage{amsmath,amssymb,mathtools}
\usepackage[scr=rsfs,cal=euler]{mathalfa}
\usepackage{xfrac}
\usepackage[unicode,hidelinks,bookmarks=false]{hyperref}
\newcommand{\ie}{i.e.\ }
\newcommand{\cf}{cf.\ }
\newcommand{\resp}{respectively\ }
\newcommand{\Subsection}{\subsection}
\providecommand{\nobreakdash}{\nobreak\hbox{-}}
\setlength{\emergencystretch}{2em}
\multlinegap0pt
\usepackage{amssymb}
\usepackage{enumerate}
\usepackage{xcolor}
\newtheorem{proposition}{Proposition}
\newtheorem{lemma}{Lemma}
\title{Schrödinger system with quintic nonlinearity: \\ spectral stability of multiple sign-changing periodic waves \\ \vspace{0.2cm}}
\author{Guilherme de Loreno, Gabriel E. Bittencourt Moraes}
\date{Source: arXiv:2601.20733v1 (CC BY 4.0)}
\begin{document}
\maketitle

\noindent\emph{Source and licence.} Schrödinger system with quintic nonlinearity: \\ spectral stability of multiple sign-changing periodic waves \\ \vspace{0.2cm}. Version arXiv:2601.20733v1 (CC BY 4.0), distributed by arXiv under the Creative Commons Attribution 4.0 International licence, http://creativecommons.org/licenses/by/4.0/, as recorded in the arXiv licence metadata for this exact version and shown on the abstract page https://arxiv.org/abs/2601.20733v1. Full licence notice: \texttt{licenses/arxiv-2601.20733-CC-BY-4.0.txt}.\par\smallskip

\noindent\textit{Editorial scope.} Counted unit: the article's lemma dif-norma (source0.tex lines 838--840). The article's complete original proof of it is delivered with it (source0.tex lines 842--897, one \texttt{proof} environment of 56 lines). No mathematics is written, repaired or re-derived: the statement and the proof are the article's own lines, in the article's own notation, and no retained line is substituted or re-flowed.  Retained uncounted as the source-local context the proof needs: lines 451--453; lines 456--475; lines 465--467; lines 469--471; lines 500--508; lines 502--504.  Nothing was omitted from any retained span, so the package carries no omission record.  The retained lines cite 2 works, whose entries are transcribed verbatim from the article's own bibliography source in the e-print and delivered with short editorial labels recorded in the item's reference mapping.  No retained line contains a figure environment, an image-inclusion command or drawing code, so no image is delivered and the package declares no asset dependency.  Editorial formatting only, in addition: an AMS \texttt{amsart} preamble that restates the article's own declarations and macros used by the retained lines, namely the environments \texttt{proposition}, \texttt{lemma} on separate counters, together with the standard packages those lines need.  The retained environments print in this document as Proposition 1, Lemma 1 in order of appearance, whereas the article prints the same items with the numbering of its own full document.  The complete native source of the article as distributed in the arXiv e-print for this exact version is bundled unchanged as \texttt{sources/193543/source0.tex}.
\begin{equation}\label{mainEDOquinticsystem}
	-\varphi''+\omega \varphi -(\kappa+\gamma B^3)\varphi^5=0,
\end{equation}

In this context, concerning the auxiliary equation
\begin{equation}\label{AuxilaryquinticEDO}
	- \phi'' + \omega \phi - \phi^5 = 0,
\end{equation}
for a periodic function $\phi: \mathbb{R} \to \mathbb{R}$ with period $L > 0$, it was shown in \cite[Section 4.3]{NataliCardosoAmaral} (see also \cite{BittencourtLoreno2022}) that for each $k \in (0,1)$, equation \eqref{AuxilaryquinticEDO} admits sign-changing periodic solutions with a cnoidal profile given by
\begin{equation}\label{cnoidalsolution}
	\phi(x) = \frac{a \, \text{cn} \left( \frac{4{\rm K}(k)}{L} x, k \right)}{\sqrt{1 - q \, \text{sn}^2 \left( \frac{4{\rm K}(k)}{L} x, k \right)}}, \quad x \in \mathbb{R}
\end{equation}
where
\begin{equation}\label{avaleuCnoidal}
	a = \frac{2}{L} \left[ \left( 2 - k^2 + 2 \sqrt{k^4 - k^2 + 1}\right) L^2 \, {\rm K}(k)^2 \right]^{\frac{1}{4}}
\end{equation}
and
\begin{equation}\label{qvaleuCnoidal}
	q = -1 + k^2 - \sqrt{k^4 - k^2 + 1}.
\end{equation}
Furthermore, the wave frequency $\omega>0$, which depends on $k \in (0,1)$ and $L>0$, is given by
\begin{eqnarray}\label{valuewCnoidal}
	\omega =  \frac{16{\rm K}(k)^2 \sqrt{k^4 - k^2 +1 }}{L^2}.
\end{eqnarray}

\begin{equation}
	a = \frac{2}{L} \left[ \left( 2 - k^2 + 2 \sqrt{k^4 - k^2 + 1}\right) L^2 \, {\rm K}(k)^2 \right]^{\frac{1}{4}}
\end{equation}

\begin{equation}
	q = -1 + k^2 - \sqrt{k^4 - k^2 + 1}.
\end{equation}

\begin{proposition}[Existence of Solutions with Cnoidal Profile]\label{cnoidalcurve}
	Let $L>0$ be fixed. Equation \eqref{mainEDOquinticsystem} admits $L$-periodic solutions with a cnoidal profile of the form
	\begin{equation}\label{cnoidal-solution}
		\varphi(x) = \frac{1}{\sqrt[4]{(\kappa + \gamma B^3)}}\frac{ a \, {\rm cn}\left( \frac{4{\rm K}(k)}{L}x,k\right)}{\sqrt{ 1 - q \, {\rm sn}^2 \left( \frac{4 {\rm K}(k)}{L} x, k \right)}},
	\end{equation}
	where the parameters $a \in \left( 2\sqrt{\tfrac{\pi}{L}}, +\infty \right)$, $q \in (-2,-1)$, and the frequency $\omega \in \left( \frac{4\pi^2}{L^2},+\infty\right)$ depend smoothly on $k\in (0,1)$ and $L>0$, as given in \eqref{avaleuCnoidal}, \eqref{qvaleuCnoidal}, and \eqref{valuewCnoidal}, respectively. In addition, the mapping
	$$\omega \in \left( \tfrac{4\pi^2}{L^2},+\infty \right) \longmapsto \varphi_\omega=\varphi \in H^2_{\rm per}$$
	is a smooth curve of periodic solutions for \eqref{mainEDOquinticsystem}.
\end{proposition}

	\begin{equation}
		\varphi(x) = \frac{1}{\sqrt[4]{(\kappa + \gamma B^3)}}\frac{ a \, {\rm cn}\left( \frac{4{\rm K}(k)}{L}x,k\right)}{\sqrt{ 1 - q \, {\rm sn}^2 \left( \frac{4 {\rm K}(k)}{L} x, k \right)}},
	\end{equation}

\begin{lemma}\label{dif-norma}
Let $L>0$ and $\omega \in \left( \frac{4\pi^2}{L^2}, +\infty \right)$ be fixed. {If  $\varphi=\varphi_\omega$ is the periodic solution with cnoidal profile given in Proposition \ref{cnoidalcurve}, then   $\frac{d}{d\omega} \|\varphi\|_{L^2}^2 > 0$.}
\end{lemma}

\begin{proof}
    Firstly, we note that $\omega$ depends smoothly on $k \in (0,1)$. By the chain rule, we have
    \begin{equation*}
        \frac{d}{d\omega} \|\varphi\|_{L^2}^2 = \frac{d}{dk} \|\varphi\|_{L^2}^2 \, \frac{dk}{d\omega} = \frac{d}{dk} \|\varphi\|_{L^2}^2 \, \left( \frac{d\omega}{dk}\right)^{-1}.
    \end{equation*}
    Our goal is to show the positivity of both $\frac{d}{dk} \|\varphi\|_{L^2}^2$ and $\frac{d\omega}{dk}$. First, using the expression for $\varphi$ given in \eqref{cnoidal-solution} and a change of variables, we obtain
    \begin{equation*}
        \|\varphi\|_{L^2}^2 = \int_0^L \varphi^2(x)\; dx = a^2 \int_0^L \frac{ \text{cn}^2 \left( \frac{4 \text{K}(k)}{L} x, k \right)}{ 1 - q \, \text{sn}^2 \left( \frac{4 \text{K}(k)}{L} x, k \right)} \; dx = \frac{a^2 L}{\text{K}(k)} \int_0^{\text{K}(k)} \frac{ \text{cn}^2(x,k)}{1-q \, \text{sn}^2(x,k)} \;dx.
    \end{equation*}
    Computing this integral via Maple and performing some simplifications, we find
    \begin{equation*}
        \|\varphi\|_{L^2}^2 = \frac{a^2 L}{\text{K}(k)} \left( \frac{\Pi(q,k) (q-1) + \text{K}(k)}{q} \right),
    \end{equation*}
    where $\Pi$ is the complete elliptic integral of the third kind defined by 
    \begin{equation*}
    \Pi(q, k) = \int_{0}^{\frac{\pi}{2}} \frac{d\theta}{(1 - q \sin^2 \theta) \sqrt{1 - k^2 \sin^2 \theta}}.
    \end{equation*}
    By substituting the expressions for $a$ and $q$ from \eqref{avaleuCnoidal} and \eqref{qvaleuCnoidal}, respectively, and employing Maple for further symbolic manipulation, we obtain
    \begin{equation}\label{ddkphi}
        \frac{d}{dk} \|\varphi\|_{L^2}^2 = \frac{\text{a}(k) \text{K}(k) + \text{b}(k) \text{E}(k)}{\text{r}(k) \sqrt{2-k^2+2\text{r}(k)} (k^2 - 1 - \text{r}(k))^2 k (1+\text{r}(k))}
    \end{equation}
    where $\text{r}(k) = \sqrt{k^4 - k^2 + 1}$, and the expressions $\text{a}(k)$ and $\text{b}(k)$ are given by
    \begin{equation*}
        \text{a}(k) = (-12k^6 + 42k^4 - 56k^2 + 32) \text{r}(k) + (12k^8 - 48k^6 + 82k^4 - 72k^2 + 32)
    \end{equation*}
    and
    \begin{equation*}
        \text{b}(k) = (-18k^4 + 40k^2 - 32) \text{r}(k) + (18k^6 - 50k^4 + 56k^2 - 32),
    \end{equation*}
    respectively. It is clear that the denominator of \eqref{ddkphi} is positive for all $k \in (0,1)$. On the other hand, the numerator $\mathcal{N}(k) = \text{a}(k) \text{K}(k) + \text{b}(k) \text{E}(k)$ satisfies
    \begin{equation*}
        \mathcal{N}(k) = \text{a}(k) \int_0^{\frac{\pi}{2}} \frac{dt}{\sqrt{ 1 - k^2 \sin^2(t)}} + \text{b}(k) \int_0^{\frac{\pi}{2}} \sqrt{ 1 - k^2 \sin^2(t)} \, dt = \int_0^{\frac{\pi}{2}} \frac{ \text{a}(k) + \text{b}(k) - \text{b}(k) k^2 \sin^2(t)}{ \sqrt{ 1 - k^2 \sin^2(t)}} \, dt,
    \end{equation*}
    which implies
    \begin{equation*}
        \mathcal{N}(k) \geq \int_0^{\frac{\pi}{2}} \left( \text{a}(k) + \text{b}(k) - \text{b}(k) k^2 \sin^2(t) \right) \, dt
    \end{equation*}
    for $t \in (0, \frac{\pi}{2})$ and $k \in (0,1)$. Thus, we obtain
    \begin{equation*}
        \mathcal{N}(k) \geq \left( \text{a}(k) + \text{b}(k) \right) \frac{\pi}{2} - \text{b}(k) k^2 \frac{\pi}{4} = \frac{\pi}{4} \underbrace{\left[ 2 \left( \text{a}(k) + \text{b}(k) \right) - \text{b}(k) k^2 \right]}_{:=\mathcal{M}(k)}.
    \end{equation*}
{Since $\mathcal{M}$ is a polynomial, it is easy to show that $\mathcal{M}$ is an increasing and positive function on $(0,1)$. Moreover, using Maple, we can verify that $\mathcal{M}(k)=0$ if and only if $k=0$. Thus, we also conclude that $\mathcal{N}(k) > 0$ for all $k \in (0,1)$, which implies $\frac{d}{dk} \|\varphi\|_{L^2}^2 > 0$ for all $k \in (0,1)$.}

    To establish that $\frac{d\omega}{dk}>0$ for all $k \in (0,1)$, we apply a similar argument. From the expression for $\omega$ in \eqref{valuewCnoidal}, we have
    \begin{equation*}
        \frac{d\omega}{dk} = \frac{16 \text{K}(k)}{L^2 (1-k^2) \, k \, \text{r}(k)} \underbrace{\left[ (2k^4 - 2k^2 + 2) \text{E}(k) + (- k^4 + 3k^2 - 2) \text{K}(k) \right]}_{:=\mathcal{P}(k)}.
    \end{equation*}
    We must show that $\mathcal{P}(k)$ is positive for all $k \in (0,1)$. Proceeding as before, we observe that
    \begin{equation*}
        \mathcal{P}(k) \geq \frac{\pi}{4} \left( 2k^4 + 2k^2 - (2k^4 - 2k^2 + 2)k^2 \right).
    \end{equation*}
{A direct verification of the polynomial on the right side of the inequality above shows that it is positive for all $k \in (0,1)$. Hence,  $\mathcal{P}(k) > 0$ for all $k \in (0,1)$, which yields $\frac{d\omega}{dk} > 0$ for all $k \in (0,1)$.  }    
%A direct verification via Maple shows that $\mathcal{P}(k) > 0$ for all $k \in (0,1)$, which yields $\frac{d\omega}{dk} > 0$ for all $k \in (0,1)$.

    Therefore, we conclude that $\frac{d}{d\omega} \|\varphi\|_{L^2}^2 > 0$ for all $k \in (0,1)$, as desired.
\end{proof}

\begin{thebibliography}{99}
\bibitem[Bitten]{BittencourtLoreno2022} Bittencourt Moraes, G. E., de Loreno, G., \textit{Cnoidal waves for the quintic Klein-Gordon and Schrödinger equations: Existence and orbital instability}, J. Math. Anal. Appl. 513 (2022) 126203. %
\bibitem[Natali]{NataliCardosoAmaral} Natali, F., Cardoso, E., Amaral, S., \textit{On the spectral stability of periodic traveling waves for the critical Korteweg-de Vries and Gardner equations}. Partial Differential Equations and Applications 2, no. 3 (2021), pp. 1-20.
\end{thebibliography}
\end{document}
